\documentstyle[%


righttag%


,11pt%


,amssymb%


,russian%               remove in Eng.


,izvru%                 Set izven% in Eng.


]{amsart}





% \text{} constructions





%\hoffset=-15mm%                  For Eng.


\setcounter{page}{3}


\god{2002}


\nomer{1~(476)} %                        Remove in Eng.


%\nomer{Vol.\,46, No.\,1}%               Set in Eng.


%\str{\pageref{begin}--\pageref{end}}%  Set in Eng.


\udk{517.98}





\author{\it Г.А.\,Акишев}


% NB:   ^^ remove in Eng.


\title{О некоторых теоремах для системы Прайса}





\date{}


%\pagestyle{myheadings}%         Set in Eng.


\pagestyle{plain} %              Remove in Eng.


\begin{document}


%\thispagestyle{plain}%          Set in Eng.


\maketitle


\label{begin}








Банахово пространство $X$ измеримых по Лебегу на отрезке $[0,1]$ функций 


называется симметричным [1], если 





$1)$ из $g\in X$ и $|f(t)|\le|g(t)|$ почти всюду на $[0,1]$ 


следует $f \in X$ и $\|f\|_X \le \|g\|_X$;





$2)$ из $f \in X$ и равноизмеримости функций $|f(t)|$ и $|g(t)|$ следует


$g \in X$ и $\|f\|_X=\|g\|_X$ ([1], с.\,123).





\noindent{}Здесь и в дальнейшем $\|\psi\|_X$ 


означает норму элемента $\psi \in X$.





Примерами симметричных пространств являются пространства Лебега $L_q$,


$1 \le q <+\infty$, Орлича, Лоренца, Марцинкевича [1].





Пусть $\chi_e(t)$ --- характеристическая функция измеримого множества 


$e \subset [0,1]$. Функция $\varphi(\mu e)=\|\chi_e\|_X$ называется 


фундаментальной функцией пространства $X$, где $\mu e$ --- мера Лебега 


измеримого множества $e \subset [0,1]$.





Таким образом, фундаментальная функция симметричного пространства $X$ 


есть функция $\varphi(t)=\|\chi_{[0,t]}\|_X$, определенная на $[0,1]$.


Функцию $\varphi(t)$ можно считать вогнутой, неубывающей, непрерывной 


на $[0,1]$, причем $\varphi(0)=0$ ([1], с.\,137).





В дальнейшем симметричное пространство с фундаментальной функцией $\varphi$ 


будем обозначать $X(\varphi)$.





Через $\sigma_{\tau}$ обозначим оператор растяжения 


$(\sigma_{\tau}f)(x)=f(\frac{x}{\tau})$, если $\frac{x}{\tau}\in [0, 1]$ и 


$(\sigma_{\tau}f)(x)=0$, если $\frac{x}{\tau}\notin [0, 1]$.





Известно, что оператор $\sigma_{\tau}$ непрерывен в любом симметричном 


пространстве и существуют пределы


$$


\underline{\gamma}_{X}=\lim_{\tau\rightarrow 0}\frac{\ln 


\|\sigma_{\tau}\|_{X\rightarrow X}}{\ln \tau},\quad 


\overline{\gamma}_{X}=\lim_{\tau\rightarrow \infty}\frac{\ln 


\|\sigma_{\tau}\|_{X\rightarrow X}}{\ln \tau},


$$


при этом $0\leq \underline{\gamma}_{X}\leq\overline{\gamma}_{X}\leq 1$ ([1], 


cс.\,131, 134; [2], с.\,1250), $\|\sigma_{\tau}\|_{X\rightarrow X}$ ---


норма оператора $\sigma_{\tau}$.





Симметричное пространство $X(\varphi)$ называется максимальным 


пространством, если 


$$


\|f\|_X=


\sup\bigg\{\int_0^1 f(t)g(t)dt :\, g \in 


X'(\overline{\varphi}),\ \|g\|_{X'} \le 1 \bigg\},


$$


где $X'(\overline{\varphi}) $ --- двойственное симметричное пространство к 


пространству $X(\varphi) $.


Отметим также, что если $\varphi(t)$ --- фундаментальная функция 


симметричного пространства $X$, то функция 


$\overline{\varphi}(t)=\frac{t}{\varphi(t)}$ при $t \in (0,1]$ и 


$\overline{\varphi}(0)=0$ является фундаментальной функцией сопряженного 


пространства $X'$ ([1], с.\,144).





Примерами максимальных симметричных пространств являются пространства Лебега 


$L_{q}$, $1\leq q<\infty$, Орлича, Лоренца, Марцинкевича ([1], с.\,144).





Для функции $\varphi(t)$, $t \in [0,1]$, положим


$$


\alpha_\varphi=\underset{t\to 0}{\varliminf}{}\frac{\varphi(2t)}{\varphi(t)} 


\quad, 


\beta_\varphi=\underset{t\to 0}{\varlimsup}{}\frac{\varphi(2t)}{\varphi(t)}.


$$


Известно, что если $X(\varphi)$ --- симметричное пространство, то 


$1\leq\alpha_{\varphi}\leq\beta_{\varphi}\leq 2.$ 


Для краткости будем пользоваться записью $A(y)\asymp B(y)$, которая


означает, что существуют положительные постоянные $C_{1}$, $C_{2}$ 


такие, что $C_{1}A(y) \leq B(y) \leq C_{2}A(y)$ для всех $y$.





Пусть дана последовательность $\{p_n\}_{n=1}^{\infty}$ целых чисел 


$p_n\geq2$, $n=1,2,\dotsc$ Положим $G=\big\{x=\{x_{n}\}_{n=1}^{\infty}: 


x_{n}$ --- целое число, $0\leq x_n\leq p_{n}-1\big\}$ и $m_{0}=1$, 


$m_{n}=p_{1}\cdot\dotsc\cdot p_{n}$.


Множество $G$ является группой с операцией 


сложения $\dotplus$, как покоординатным сложением по модулю $p_{n}$, 


$n=1,2,\dotsc$ Топология в группе $G$ определяется системой подгрупп


$$


G_{n}=\big\{x=\{x_{k}\}_{k=1}^{\infty}\in G:


x_{k}=0 \text{ для } k=1,2,\dotsc,n\big\}.


$$


 


Пусть отображение $\lambda:G\rightarrow[0,1]$ определено равенством


$$


\lambda(\wt{x})=\sum_{k=1}^{\infty}\frac{x_{k}}{m_{k}}=x,


$$


где $\wt{x}\in G$ и $x\in [0,1]$, $x_{k}=0,1,\dotsc,p_{k}-1$.


Это отображение является взаимно однозначным всюду, кроме точек 


$\frac{l}{m_{n}}$, $l=0, 1,\dotsc, m_{n}-1$, $n=1,2,\dotsc$ Заменяя отрезок 


$I\equiv[0,1]$ модифицированным отрезком $I^{*}$ с соответствующими 


топологией и операцией сложения, как это сделано в [3], получим изоморфные 


топологические группы $I^{*}$, $G$.





На модифицированном отрезке $I^{*}=[0,1]^{*}$ определим систему функций 


Прайса [3]. Положим 


$$


\psi_{0}(x)=1,\quad \psi_{m_{k}}(x)=\exp


\bigg(\frac{2\pi i x_{k+1}}{p_{k+1}}\bigg),\ \ x\in I^{*}.


$$


Если $n=\sum\limits_{k=0}^{r}\alpha_{k}m_{k}$, где


$\alpha_{k}=0,1,\dotsc, p_{k+1}-1$, то положим


$$


\psi_{n}(x)=\prod_{k=0}^{r}(\psi_{m_{k}}(x))^{\alpha_{k}},\ \ x\in I^{*}.


$$


Система функций $\{\psi_{n}\}$ является полной ортонормированной периодической


мультипликативной системой [3].





Через $a_{n}(f)$ будем обозначать коэффициенты Фурье 


функций $f\in L_{1}$ по системе


Прайса $\{\psi_{n}\}$. Положим


$$


\sigma_{n,j}(f,t)=\sum_{v=jm_{n}}^{(j+1)m_{n}-1}a_{v}(f)\psi_{v}(t),


\ \ j=0,1,\dotsc,p_{n+1}-1,\ \ t\in[0,1].


$$


В случае $\sup\limits_{n}p_{n}<+\infty$ Ватари [4] доказал, что для любой 


функции $f\in L_{q}$, $1<q<+\infty$,


имеют место неравенства


\begin{equation}


c_{1}\|f\|_{q}\leq\bigg\|\bigg(|a_{0}(f)\mid^{2}+\sum_{n=0}^{\infty} 


\sum_{j=0}^{p_{n+1}-1}|\sigma_{n,j}(f,t)|^{2}\bigg)^{\frac{1}{2}}\bigg\|_{q}


\leq c_{2}\|f\|_{q},


\end{equation}


где $c_{1}$, $c_{2}$ --- некоторые положительные постоянные. Здесь $L_{q}$ 


--- пространство Лебега с нормой 


$$


\|f\|_{q}=\bigg(\int_{0}^{1}|f(x)|^{q}dx\bigg)^{\frac{1}{q}},\quad 1\leq 


q<+\infty .


$$


Однако в случае $\sup\limits_{n}p_{n}=+\infty$ неравенства (1) не всегда 


верны [4], [5].





В данной статье эти вопросы изучаются для функции $f\in X(\varphi)$.


Для доказательства основных результатов используется


\medskip





{\bf Лемма A} ([6], с.\,10). {\it 


Пусть даны $\psi$-функции


$\varphi_{1}(x)$ и $\varphi_{2}(x)$, $x\in [0,1]$ и


$\beta_{\varphi_{1}}<\alpha_{\varphi_{2}}$. Тогда для функции


$$


\theta (x)=


\begin{cases}


\frac{\varphi_{2}(x)}{\varphi_{1}(x)},& x\in (0, 1];\\


0,& x=0,


\end{cases}


$$


существует $\psi$-функция $\theta_{1}(x)$, для которой 


$\alpha_{\theta_{1}}>1$ и $\theta_{1}(x)\asymp\theta(x)$, $x\in [0, 1]$. 


}





\begin{thm} Пусть $X(\varphi)$ --- симметричное пространство, 


$0<\underline{\gamma}_{X},\, \overline{\gamma}_{X}<1$ и 


мультипликативная система $\{\psi_{n}\}$ определена ограниченной 


последовательностью $\{p_{n}\}$. Тогда для любой функции $f\in 


X(\varphi)$ имеют место неравенства


$$


c_{1}\|f\|_{X}\leq\bigg\|\bigg(| 


a_{0}(f)|^{2}+\sum_{n=0}^{\infty} \sum_{j=0}^{p_{n+1}-1} 


|\sigma_{n,j}(f)|^{2}\bigg)^{\frac{1}{2}}\bigg\|_{X}\leq 


c_{2}\|f\|_{X}.


$$


\end{thm}





\begin{pf}


Так как по условию теоремы $0<\underline{\gamma}_{X}\leq 


\overline{\gamma}_{X}<1$, то существуют $q_{0}, q_{1}\in (1, \infty)$ такие, 


что $0<\frac{1}{q_{0}}< \underline{\gamma}_{X}\leq 


\overline{\gamma}_{X}<\frac{1}{q_{1}}<1$. 


В силу правого неравенства в (1) сублинейный оператор 


$$


Tf=\bigg(| a_{0}(f)|^{2}+\sum_{n=0}^{\infty} \sum_{j=0}^{p_{n+1}-1}


|\sigma_{n,j}(f)|^{2}\bigg)^{\frac{1}{2}}


$$


ограниченно действует из $L_{q_{i}}$ в $L_{q_{i}}$ ($i=0,1$). Следовательно,


этот оператор ограниченно действует из $X(\varphi)$ в $X(\varphi)$ [7], 


[8], т.\,е.


$$


\bigg\|\bigg(| a_{0}(f)|^{2}+\sum_{n=0}^{\infty} \sum_{j=0}^{p_{n+1}-1}


|\sigma_{n,j}(f)|^{2}\bigg)^{\frac{1}{2}}\bigg\|_{X}\leq 


c\|f\|_{X}


$$


для любой функции $f\in X(\varphi)$.





Противоположное неравенство доказывается с помощью принципа двойственности и


предыдущего неравенства.


\end{pf}





\begin{thm}


Пусть $X(\varphi)$ --- максимальное симметричное 


пространство, $2^{\frac{1}{2}}<\alpha_{\varphi}\leq\beta_{\varphi}<2$, и 


система Прайса $\{\psi_{n}\}$ определена неограниченной 


последовательностью $\{p_{n}\}$. Тогда существует функция 


$f_{0}\in X(\varphi)$, для которой


$$ 


\sum_{n=0}^{\infty} 


\sum_{j=1}^{p_{n+1}-1}|\sigma_{n,j}(f_{0},t)|^{2}=+\infty


$$


для всех $t\in[0,1]$. 


\end{thm}





\begin{pf}


Так как последовательность $\{p_{n}\}$ неограниченная, 


то можно выбрать подпоследовательность


$\{p_{n_{k}}\}$ такую, что $p_{n_{k}+1}>2^{k}$, $k=1,2,\dotsc$


Рассмотрим ряд


$$


\sum_{n=0}^{\infty} \sum_{j=1}^{p_{n+1}-1}\sigma_{n,j}(t),\ \ t\in[0,1],


$$


с коэффициентами


$a_{v}=(k p_{n_{k}+1})^{-\frac{1}{2}}$, если $v=j m_{n_{k}}$,


$a_{v}=0$, если $v\neq j m_{n_{k}}$.


Покажем, что этот ряд по норме пространства $X(\varphi)$ сходится к 


некоторой функции $f_{0}\in X(\varphi)$. По свойству нормы будем иметь


\begin{equation}


\bigg\|\sum_{k=l}^{v} \sum_{j=1}^{p_{n_{k}+1}-1}\sigma_{n_{k},j}(\cdot) 


\bigg\|_{X}\leq\sum_{k=l}^{v} 


\bigg\|\sum_{j=1}^{p_{n_{k}+1}-1}\sigma_{n_{k},j}(\cdot) \bigg\|_{X}.


\end{equation}





Пусть $g\in X'(\overline{\varphi})$ --- произвольная функция такая, что 


$\|g\|_{X'}\leq 1$. Тогда по определению функций системы Прайса получим


\begin{multline}


\bigg|\int_{0}^{1}\ \sum_{j=1}^{p_{n_{k}+1}-1}


\sigma_{n_{k},j}(t)g(t)dt\bigg|=


\bigg|\int_{0}^{1}\ \sum_{j=1}^{p_{n_{k}+1}-1}(k p_{n_{k}+1})^{-\frac{1}{2}} 


\psi_{j m_{n_{k}}}(t)g(t)dt\bigg|= \\


%


=(k p_{n_{k}+1})^{-\frac{1}{2}}\bigg|\int_{0}^{1}


\ \sum_{j=1}^{p_{n_{k}+1}-1}\bigg(X_{E(k)}(t)+X_{\ov{E}(k)}(t)


\exp\bigg\{2\pi i\frac{j_{n_{k}+1}(t)}{p_{n_{k}+1}}\bigg\}


\bigg)g(t)dt\bigg|.


\end{multline}


Здесь $X_{E(k)}(t)$, $X_{\ov{E}(k)}(t)$ --- характеристические функции 


соответственно множеств


$E(k)=\bigcup\limits_{r=0}^{m_{n_{k}-1}}


[\frac{r}{m_{n_{k}}},\frac{r}{m_{n_{k}}}+\frac{1}{m_{n_{k}+1}}]$


и


$\ov{E}(k)=[0,1]\setminus E(k)$.





По определению максимального симметричного пространства и 


характеристической функции, учитывая, что мера Лебега


$\mu E(k)=\frac{1}{p_{n_{k}+1}}$, получим


\begin{multline*}


\bigg|\int_{0}^{1}\ \sum_{j=1}^{p_{n_{k}+1}-1}\bigg(X_{E(k)}(t)+


X_{\ov{E}(k)}(t) \exp\bigg\{2\pi i\frac{j_{n_{k}+1}(t)}


{p_{n_{k}+1}}\bigg\}\bigg)g(t)dt\bigg|= \\


%


=\bigg|(p_{n_{k}+1}-1)


\int_{E(k)}g(t)dt-\int_{\ov{E}(k)}g(t)dt\bigg|\leq


(p_{n_{k}+1}-1)\|X_{E(k)}\|_{X}+\|X_{\ov{E}(k)}\|_{X}= \\


%


=(p_{n_{k}+1}-1)\varphi(\mu E(k))+


\|X_{\ov{E}(k)}\|_{X}=(p_{n_{k}+1}-1)\varphi(p_{n_{k}+1}^{-1})+


\|X_{\ov{E}(k)}\|_{X} 


\end{multline*}


для любой функции $g\in X'(\overline{\varphi})$, $\|g\|_{X'}\leq 1$,


откуда с учетом (3) следует


\begin{equation}


\bigg\| \sum_{j=1}^{p_{n_{k}+1}-1}\sigma_{n_{k},j}(\cdot) \bigg\|_{X}\leq


\left(\sqrt{(kp_{n_{k}+1})}\,\right)^{-1}(p_{n_{k}+1}-1)


\varphi(p_{n_{k}+1}^{-1}) +\|X_{[0,1]}\|_{X}


\end{equation}


для любого $k=1,2,\dotsc$


По условию теоремы $\sqrt{2}<\alpha_{\varphi}$. Тогда по лемме A существует


$\psi$-функ\-ция $\theta{(t)}$ такая, что


$$


\frac{\varphi(t)}{\sqrt{t}}\asymp \theta(t), \ \ t\in(0,1] 


\ \text{ и }\ \alpha_{\theta}>1.


$$


По свойству монотонности функции $\theta(t)$, учитывая, что $p_{n_k+1}>2^k$,


$k=1,2,\dotsc$, получим


\begin{equation}


\sum_{k=l}^{v}\varphi(p_{n_{k}+1}^{-1})\leq 


\sum_{k=l}^{v}\theta\bigg(\frac{1}{2^k}\bigg)\leq


(\ln2)^{-1}\int\limits_{0}^{2^{-l+1}} 


\theta(t)\frac{dt}{t}=\underline{O}(\theta(2^{-l+1})).


\end{equation}


Так как $\theta(t)\rightarrow 0$ при $t\rightarrow 0$, то из (2), (4) и (5) 


следует


$$


\bigg\|\sum_{k=l}^{v}\sum_{j=1}^{p_{n_k+1}-1}\sigma_{n_{k}j}


\bigg\|_{X}\rightarrow 0,\quad v, l\rightarrow +\infty.


$$


Следовательно, в силу полноты пространства $X(\varphi)$ существует 


функция $f_0\in X(\varphi)$ такая, что


$$


\bigg\|f_0-\sum_{k=1}^{v}\sum_{j=1}^{p_{n_k+1}-1}\sigma_{n_{k}j}\bigg\|_{X}


\rightarrow 0,\quad v\rightarrow +\infty,


$$


т.\,е. числа $a_v$ будут коэффициентами Фурье функции $f_0\in X(\varphi)$ по 


мультипликативной системе \{$\psi_n$\}.


Но легко видеть, что


$$


\sum_{k=0}^{\infty}\sum_{j=1}^{p_{n_k+1}-1}|\sigma_{n,j}(t)|^{2}=+\infty


$$


для любого $t\in[0,1]$.


\end{pf}





\begin{thm}


Пусть система Прайса определена неограниченной 


последовательностью \{$p_n$\}. Тогда существует функция $f(t)$, не 


принадлежащая ни одному из максимальных симметричных пространств 


$X(\varphi)$, $1<\beta_{\varphi}<\sqrt{2}$, для которой


$$


\sum_{k=0}^{\infty}\sum_{j=1}^{p_{n_k+1}-1}|\sigma_{nj}(f,t)|^{2}\leq c


$$


для любого $t\in[0,1]$.


\end{thm}





\begin{pf}


Пусть \{$p_{n_k}$\} --- подпоследовательность, построенная при 


доказательстве теоремы 2. Так как по условию теоремы 


$\beta_{\varphi}<\sqrt{2}$, то существует $\varepsilon \in (0,\frac{1}{2})$ 


такое, что $\beta_{\varphi}<2^{\frac{1}{2}-\varepsilon}<\sqrt{2}$.





Рассмотрим ряд


$$


\sum_{n=0}^{\infty}\sum_{j=1}^{p_{n+1}-1}\sigma_{n,j}(t), \quad t\in[0,1],


$$


с коэффициентами


$$


a_v=


\begin{cases}


p_{n_k+1}^{-\frac{1}{2}-\varepsilon}, & \text{ если } v=jm_{n_{k}};\\


0, & \text{ если } v\neq jm_{n_{k}}.


\end{cases}


$$


Тогда легко убедиться, что


$$


\sum_{n=0}^{\infty}\sum_{j=1}^{p_{n+1}-1}|\sigma_{n,j}(t)|^2=


\sum_{k=1}^{\infty}\sum_{j=1}^{p_{n_k+1}-1}|a_{jm_{n_k}}|^2=


\sum_{k=1}^{\infty}p_{n_k+1}^{-2\varepsilon}<+\infty


$$


для любого $t\in[0,1]$. 


Этим одновременно доказали, что данный ряд будет рядом Фурье


некоторой функции $f\in L_2$.





Теперь покажем, что функция $f\in L_2$ 


не принадлежит ни одному из пространств


$X(\varphi)$, $\beta_{\varphi}<\sqrt{2}$. По определению максимального 


симметричного пространства $X(\varphi)$ и функций Прайса получим


\begin{equation}


\bigg\|\sum_{j=1}^{p_{n_k+1}-1}a_{jm_{n_k}}\psi_{jm_{n_k}}\bigg\|_{X} \geq 


\frac{1}{\|X_{E(k)}\|_{X'}} \bigg|\int_{E(k)}\sum_{j=1}^{p_{n_k+1}-1}


a_{jm_{n_k}}\psi_{jm_{n_k}}(t)dt\bigg|=


\frac{1}{2} p_{n_k+1}^{\frac{1}{2}-\varepsilon}


\varphi(p_{n_k+1}^{-1}).


\end{equation}


По лемме A существует функция $\theta(t)$ такая, что 


$$


\frac{t^{\frac{1}{2}-\varepsilon}}{\varphi(t)}\asymp\theta(t)\rightarrow 


0\quad \text{при}\quad t\rightarrow 0.


$$


Следовательно, в силу (6)


$$


\bigg\|\sum_{j=1}^{p_{n_k+1}-1}a_{jm_{n_k}}\psi_{jm_{n_k}}\bigg\|_{X}


\rightarrow +\infty \quad \text{при} \quad k\rightarrow+\infty.


$$





Значит, $f\notin X(\varphi)$.


\end{pf}





{\bf Замечание.} В случае $X(\varphi)=L_{q}$, $1<q<+\infty$, $q\neq 2$, 


теоремы 2 и 3 ранее доказаны в [5].


Для системы Уолша (т.\,е. $p_n=2$ для всех $n=1,2,\dotsc$) неравенства (1) в 


пространстве Орлича доказаны в [9].








\begin{thebibliography}{9}





\bibitem{1}


Крейн С.Г., Петунин Ю.И., Семенов Е.М. {\em Интерполяция линейных 


операторов}. -- М.: Наука, 1978. -- 400 с.





\bibitem{2}


Boyd D.W. {\em Indices of function spaces and their relationship to 


interpolation} // Can. J. Math. -- 1969. -- V.\,21.


-- \No\,5. -- P.\,1245--1254.





\bibitem{3} 


Голубов Б.И., Ефимов А.В., Скворцов В.А. {\em Ряды и преобразования


Уолша$:$ Теория и применения}. -- М.: Наука, 1987. -- 344 с.





\bibitem{4} 


Watari C. {\em On generalized Walsh Fourier series} // T\^o{}hoku %% place --


Math. J. -- 1958. -- V.\,10. -- \No\,3. -- P.\,211--241.





\bibitem{5} 


Vlasova E.A. {\em Convergence of series with respect to generalized


Haar systems} // Anal. Math. -- 1987. -- V.\,13. -- \No\,4.


-- P.\,339--360.





\bibitem{6} 


Лапин С.В. {\em Некоторые теоремы вложения для произведений функций}. 


-- Моск. гос. ун-т. -- М., 1980. -- 31 с. -- Деп. в ВИНИТИ 18.03.80,


\No\,1036-80. 





\bibitem{7} 


Janson S. {\em On the interpolation of sublinear operators} // 


Studia Math. -- 1982. -- V.\,75. -- P.\,51--53.





\bibitem{8} 


Бухвалов А.В. {\em Интерполяция операторов в пространствах вектор-функций с 


приложениями к сингулярным интегральным операторам} // ДАН СССР.


-- 1984. -- T.\,278. -- \No\,3. -- C.\,523--526.





\bibitem{9} 


Yung-min Chen. {\em Theorems of asymptotic approximation} // 


Math. Ann. -- 1960. -- V.\,140. -- P.\,360--407.





\end{thebibliography}





\vskip 2cm





\postup{\it Карагандинский государственный }{\qquad \it Поступили}


\postup{\it университет им. Е.А.\,Букетова}{\it первый вариант $10.01.2000$}


\postup{}{\it окончательный вариант $11.04.2001$}





\label{end}


\end{document}









































